\subsection{初步推论}

C6. 设 \(A, B, C\) 是 \(\mathfrak{C}\) 中的关系。如果 \(A \Longrightarrow B\) 并且 \(B \Longrightarrow C\) 是定理 \(\mathfrak{C}\) ，那么 \(A \Longrightarrow C\) 是 \(\mathfrak{C}\) 中的定理 .

事实上，\((\pmb{B} \Rightarrow \pmb{C}) \Rightarrow ((\pmb{A} \Rightarrow \pmb{B}) \Rightarrow (\pmb{A} \Rightarrow \pmb{C}))\) 是 \(\mathfrak{C}\) 的一个公理；这是在S4中把 \(\pmb{A}\) 换成 \(\pmb{B}\) 、把 \(\pmb{B}\) 换成 \(\pmb{C}\) ，并把 \(\pmb{C}\) 换成“非 \(\pmb{A}\)”得到的。根据C1（§2，第2节），\((\pmb{A} \Rightarrow \pmb{B}) \Rightarrow (\pmb{A} \Rightarrow \pmb{C})\) 是 \(\mathfrak{C}\) 中的定理。再应用一次C1即完成证明。

C7. 如果A和B是C中的关系，则 \(\mathbf{B} \Longrightarrow (\mathbf{A} \text{ 或 } \mathbf{B})\) 是C中的一个定理。

因为 \(\boldsymbol{B}\Longrightarrow(\boldsymbol{B}\) 或 A）和 \((\boldsymbol{B}\) 或 A） \(\Longrightarrow(\boldsymbol{A}\) 或 B）根据S2和S3都是C的公理。现在应用C6。

C8. 如果 \(A\) 是 \(\mathfrak{C}\) 中的一个关系， \(A \Longrightarrow A\) 是 \(\mathfrak{C}\) 中的定理 .

因为 \(A \Longrightarrow (A \text{ 或 } A)\) 和 \((A \text{ 或 } A) \Longrightarrow A\) 根据S2和S1都是公理。现在应用C6。

C9. 如果A是关系且B是T中的定理，则 \(A \Longrightarrow B\) 是T中的定理。

因为根据C7，\(\boldsymbol{B} \Rightarrow ((\text{非 } \boldsymbol{A}) \text{ 或 } \boldsymbol{B})\) 是一个定理，所以“（非 \(\boldsymbol{A}\) ）或 \(\boldsymbol{B}\) ”是定理；换言之，根据C1，\(\boldsymbol{A} \Rightarrow \boldsymbol{B}\) 是定理。

C10. 如果A是T中的关系，则“A或(非A)”是T中的定理。

因为“(非A)或A”根据C8是定理；现在使用S3和C1。

C11. 若 A 是 C 中的一个关系，则“A \(\Longrightarrow\) （非非A）”是 C 中的一个定理。

因为这个关系就是“（非A）或（非非A）”，结果由C10得出。

C12. 设A和B是C中的两个关系。则关系

\[
(A \Rightarrow B) \Rightarrow ((\text { 非 } B) \Rightarrow (\text { 非 } A))
\]

是 \(\mathfrak{C}\) 中的定理 .

因为

\[
((\text { 非 } A) \text { 或 } B) \Longrightarrow ((\text { 非 } A) \text { 或 } (\text { 非   非 } B))
\]

是一个定理，根据C11、S4和C1。另一方面，

\[
((\text { 非 } A) \text { 或 } (\text { 非   非 } B)) \Longrightarrow ((\text { 非   非 } B) \text { 或 } (\text { 非 } A))
\]

是公理，根据S3。因此

\[
((\text { 非 } A) \text { 或 } B) \Longrightarrow ((\text { 非   非 } B) \text { 或 } (\text { 非 } A))
\]

由C6可知是一个定理。因此结论成立。

C13. 设 \(A, B, C\) 是 \(\mathfrak{C}\) 中的关系。如果 \(A \Rightarrow B\) 是 \(\mathfrak{C}\) 中的定理 ，那么 \((B \Rightarrow C) \Rightarrow (A \Rightarrow C)\) 是 \(\mathfrak{C}\) 中的定理 .

因为（非B） \(\Longrightarrow\) （非A）是一个定理，依据C12和C1。因此（C或（非B）） \(\Longrightarrow\) （C或（非A））是一个定理，依据S4和C1。通过两次应用S3和C6，我们推断出

\[
((\text { 非 } \boldsymbol {B}) \text { 或 } \boldsymbol {C}) \Longrightarrow ((\text { 非 } \boldsymbol {A}) \text { 或 } \boldsymbol {C})
\]

是一个定理；但这正是给定的关系。

¶ 从现在起，我们通常将不加明确引用地使用C1和C6。

